\section{Planning}



\begin{frame}{Bias for Action}
    
    \mytext{
    \begin{itemize}
        \item Tendency to act without plan
        \item Makes sense when there are no risks
        \item Example: business-as-usual
        \item Planning often under-financed or not financed!?
        \item Optimism bias
    \end{itemize}
    }
\end{frame}


\begin{frame}{Agenda, to-do-lists}
    
    \mytext{
    \begin{itemize}
        \item \emph{agenda} n.pl.~lat.~`(things) to be done'
        \item problem: top priority may delay everything else
        \item things stay on list until they become urgent
        \item possible multi-tasking?
        \item division of labour?
        \item better: define a date (point in time)! $\rightarrow$ \textit{schedule!}
        \item have multiple agendae for different dates
        \item colour scheme (red-yellow-green)
        \item \textbf{software}
    \end{itemize}
    }
\end{frame}


\begin{frame}{'Eisenhower' Matrix}
    \showmytikz{../shared/tikz/tikz-eisenhower}
    \mycomment{
        \vfill
    Source: own design / commons
    }
\end{frame}


\begin{frame}{Project Management Software}
    
    \mytext{
    \begin{itemize}
        \item Why?
        \item What to look for?
        \item Criteria?
    \end{itemize}
    }
\end{frame}


\begin{frame}{Foresight, Planning}
    
    \mytext{
    \begin{itemize}
        \item Efficiency comes with planning. 
        \item \textbf{Planning dominates the management of complex  projects.}
        \item Complexity may not scale linearly.
        \item Planning is a communicative activity,
        \item Planning has to do with foresight
        \item \textit{What must be done?}
        \item \textit{And in which order?}
        \item \textit{By whom, with what?}
        \item \textit{Depending on what?}
        \item What are the sub-projects or tasks?
        \item \textbf{Plans change}
    \end{itemize}
    }
\end{frame}

% \begin{frame}{Objective and Owner of a Project}
%     \begin{itemize}
%         \item \textbf{Who }wants what?
%         \item Project owner
%         \item Project sponsor
%         \item Project governance
%         \item Common vision?
%         \item Well defined?
%     \end{itemize}
% \end{frame}

% \begin{frame}{Project Flow Chart (NEW)}
%     \showmytikz{../shared/tikz/tikz-projectflowchartfour}
    
% \end{frame}

\begin{frame}{Project Flow Chart}
    
    \mytext{
    \begin{itemize}
        \item X initialised the project
        \item to achieve Y;
        \item Business environment = project began
        \begin{itemize}
            \item planning
            \item execution
            \item closure
        \end{itemize}
        \item X was achieved / running
    \end{itemize}
    }
\end{frame}

\begin{frame}{Project Flow Chart}    
    \showmytikz{../shared/tikz/tikz-PM-project-flow-linear}
    \end{frame}

\begin{frame}{Project Flow Chart}
    \showmytikz{../shared/tikz/tikz-PM-project-flow-branching}
    \end{frame}

\begin{frame}{Project Flow Chart}    
    \showmytikz{../shared/tikz/tikz-PM-project-flow-branching-labeled}
    
    
\end{frame}

\begin{frame}{What needs to be done?}
    
    \mytext{
    \begin{itemize}
        \item \textbf{Planning starts with listing tasks}
        \item[] \textbf{and defining their deliverables}
        \item[] 
        \item \textbf{Planning is scheduling tasks}
        \item[] \textbf{and allocating resources}
        \item[] 
        \item \textbf{Plans change!}
    \end{itemize}
    }
\end{frame}

\begin{frame}{Sub-Projects are called 'Work Packages'}
    
    \mytext{
    
    Sub-projects are often distributed
    
    \medskip
    
    \textbf{`Work package' (WP)}
    
    \medskip
    
    German: `Arbeitspaket (AP)' - `Teilarbeitspaket'\\
    English: Subproject, Task, Work Package
    }
\end{frame}



\begin{frame}{Break Work down to Tasks}

    \begin{forest}
        for tree={
            font=\footnotesize,
            grow = east,
            anchor=west,
              forked edge,
            inner sep=2pt
                }
        [Project
            [Work package ... [...][...]]
            [Work package 3 [Task 3.2 ][Task 3.1 ]]
            [Work package 2 [Task 2.2 ][Task 2.1 ]]
            [Work package 1 [Task 1.3 ][Task 1.2 ][Task 1.1 ]]
        ]
    \end{forest}
    \end{frame}

\begin{frame}{Break Work down to Tasks}    
        \begin{forest}
        for tree={
            font=\footnotesize,
            grow = east,
            anchor=west,
            forked edge,
            inner sep=2pt
                }
        [Project
            [Work package ... [...][...]]
            [Work package 3 [Task 3.2 - Alice][Task 3.1 - Alice]]
            [Work package 2 [Task 2.2 - Charlie][Task 2.1 - Bob]]
            [Work package 1 [Task 1.3 - Doreen][Task 1.2 - Alice][Task 1.1 - Doreen]]
        ]
    \end{forest}
    \end{frame}

\begin{frame}{Break Work down to Tasks}    
        \begin{forest}
        for tree={
            font=\footnotesize,
            grow = east,
            anchor=west,
            forked edge,
            inner sep=2pt
                }
        [Project
            [Work package ... [...][...]]
            [Work package 3 [Task 3.2 - Alice - 12d][Task 3.1 - \textit{Alice - 15d}]]
            [Work package 2 [Task 2.2 - Charlie - 6d][Task 2.1 - Bob - 8d]]
            [Work package 1 [Task 1.3 - Doreen - 9d][Task 1.2 - Alice - 4d][Task 1.1 - Doreen - 9d]]
        ]
    \end{forest}
    
\end{frame}


% \input{TE_Hatfactory}

\begin{frame}{Flow Charts, Lists, Tree Chart, Gantt Diagramm --}
    
    \mytext{
You need an overview of what is to be done!

The overview can take the form

\medskip

    \begin{itemize}
        \item \colorbox{green}{flow chart},a.k.a. \textit{network diagram}
        \item \colorbox{lightgray}{List}, a.k.a. \textit{register}
        \item \colorbox{yellow}{WBS}, a.k.a. \textit{tree diagram}
        \item \colorbox{red}{Gantt Diagram}
    \end{itemize}
    }
\end{frame}

\begin{frame}{Flow Charts, Lists, Tree Chart, Gantt Diagramm --}
    \showmytikz{../shared/tikz/tikz-wbs-tree-gantt-list}
    \mycomment{
    \textbf{You always need all four forms!}
    
    \medskip
    
    

        
    }
\end{frame}


\begin{frame}{Project "Kitchen Renovation"}
    \showmytikz{../shared/tikz/tikz-chain-kitchen}
    \mycomment{
    \colorbox{green}{\textbf{Flowchart}}
    
    \medskip
    
    
    }
\end{frame}

\begin{frame}{Project "Kitchen Renovation"}
    \resizebox{\textwidth}{!}{
        \begin{tabular}{rllrrllr}
        \toprule
        Task & Process & Result & hrs. & Dep. & Assigned & Resources & Budget\\
        \midrule
        1. & Manage & (the project) & cont. & management & Alice & Telephone & 100 \\
        2. & Empty room & room empty & 1 & movers & Bob & hand truck, space & 100\\
        3. & Paint room & paint dry & 2 & painters & Charlie & paint etc.&  100 \\
        4. & Change floor & new floor & 2 & floorlayers & Bob & floorboards & 300 \\
        5. & Clear out & rubbish disposed & 2 & movers & Bob & van & 100 \\
        6. & Replace items & new appliancies & 2 & management & Alice & suppliers/consultants  & 4000 \\
        7. & Refurbish & livable room & 2 & Users? & Alice & movers & 500 \\
  
              \bottomrule
        \end{tabular}
        }
    \mytext{
    
    \colorbox{lightgray}{\textbf{List / register}}
    
    \medskip
    
    
        
        
    }
\end{frame}

\begin{frame}{Project "Kitchen Renovation"}
    \showmytikz{../shared/tikz/forest-houshold-wbs}
    \mycomment{
    
    \colorbox{yellow}{\textbf{WBS Work Breakdown Structure}}
    
    \medskip
    
    
    
    }
\end{frame}

\begin{frame}{Project "Kitchen Renovation"}
    \showmytikz{../shared/tikz/gantt-kitchen}\showmytikz{../shared/tikz/tikz-PM-gantt-kitchen-renovation}
    \mycomment{
    
    \colorbox{red}{\textbf{Gantt-Diagram}}
    
    \medskip
    
    % 
                
    %% is not the right one, took another one. Make new?
    
        
    }
\end{frame}


% \input{TE_Hatfactory}



% \begin{frame}{Example for Work Breakdown Structure}
%   \begin{forest}
%     for tree={
%         font=\footnotesize,
%         grow = east,
%         anchor=west,
%         forked edge,
%         inner sep=2pt
%             }
%     [\textbf{Have a new ship}
%     ]
%     \end{forest}
    
%     \framebreak
    
%     \begin{forest}
%     for tree={
%         font=\footnotesize,
%         grow = east,
%         anchor=west,
%         forked edge,
%         inner sep=2pt
%             }
%     [\textbf{Have a new ship}
%         [SP3 use ship]
%         [SP2 order ship]
%         [SP1 plan ship]
%     ]
%     \end{forest}    
    
%     \framebreak
    
%     \begin{forest}
%     for tree={
%         font=\footnotesize,
%         grow = east,
%         anchor = west,
%         forked edge,
%         inner sep=2pt
%             }
%     [\textbf{Have a new ship}
%         [SP3 use ship [find port][find jobs for ship][find crew]]
%         [SP2 order ship [supervise][contract][find contractor]]
%         [SP1 plan ship [technical planning][business planning]]
%     ]
%     \end{forest}        

% \end{frame}


\begin{frame}{Data Structure of Tasks}

    \begin{tabular}{l l l l l l}
    \toprule
         Task & Result &  Resp. & Duration & Budget & Depend. \\
         \midrule
         1 & 1.1 & Doreen & 9d & 5K & -\\
         2 & 2.1 & Alice & 4d & 2K & 1.1 \\
           & ... & ... & ... & ... & ... \\
         \bottomrule
    \end{tabular}
    \end{frame}

\begin{frame}{Data Structure of Tasks}
    
    \mytext{
    
    Every task is an entry in a database with...
    
     \begin{itemize}
         \item Area, or 'Work package'
         \item Task number
         \item Task name
         \item Duration
         \item Result, or Milestone
         \item Dependencies (backward, forward)
         \item Responsibility
         \item Budget
         \item Other ressources
     \end{itemize}
    }
\end{frame}

\begin{frame}{Data Structure of Tasks}
    
    \mytext{
    Every task is an entry in a database with...
    
     \begin{itemize}
         \item Area, or 'Work package'
         \item Task number
         \item Task name
         \item Duration
         \item Result, or Milestone
         \item Dependencies (backward, forward)
         \begin{itemize}
             \item Responsibility
             \item Budget
             \item Other ressources
         \end{itemize}
     \end{itemize}
         
    
    }
\end{frame}

\begin{frame}{Data Structure of Tasks}
    
    \mytext{
    
     IMPORTANT: 
     
     \begin{itemize}
         \item Every \textbf{task} or \textbf{work} is a \textbf{process!}
         \item[] \textit{in contrast}
        \item every \textbf{result} is a \textbf{milestone!} 
        \begin{itemize}
            \item a milestone can only be what is a result!
            \item processes are not results and not milestones!
            \item milestones are not processes!
        \end{itemize}
        \item a (result $\equiv$ milstone) can be checked!
        \item a (result $\equiv$ milstone) is an interface!
        \item a (result $\equiv$ milstone) leads to another task!
        \item all required resources are themselves dependencies and could in theory be modelled as tasks as well!
        \item however, cost, man hours, duration etc. are KPIs
        \item[] (Key Performance Indicators).
     \end{itemize}
    }
\end{frame}




\begin{frame}{How to create a WBS}
    
    \mytext{
   
        First, with expert knowledge...
   
        \begin{enumerate}
            \item brainstorm what needs to be done!
            \item group these tasks
            \item[$\rightarrow$] flow chart, tree or mindmap format
        \end{enumerate}
        
    }
\end{frame}

\begin{frame}{How to create a WBS}
    
    \mytext{
        
        Then,for each task...
        
        \begin{enumerate}
            \item define goal ($\equiv$ result, milestone)
            \item determine duration
            \item determine requirements
            \begin{itemize}
                \item dependency on other milestones
                \item resources required
            \end{itemize}
            \item assign responsibility
            \begin{itemize}
                \item identify required competency
            \end{itemize}
        \end{enumerate}
    }
\end{frame}

\begin{frame}{Flow Charts, Lists, Tree Chart, Gannt Diagramm}
    
    \showmytikz{../shared/tikz/tikz-wbs-tree-gantt-list}
        
\end{frame}

\begin{frame}{Ideal PM Software Solution}
    
    \mytext{
    
    
    We seek a software that
    
    \begin{enumerate}
        \item Allows us to mindmap WBS
        \item turns this into a database (list)
        \item shows this as Gantt-Diagram
        \item allows us to move things 
        \item allows us to record changes
        \item allows us to monitor progress
        \item allows us to monitor costs
        \item allows the team to use it
        \item allows export into presentable forms
    \end{enumerate}
    
    }
\end{frame}



\begingroup
    \ifdefined\setbeamercolor\setbeamercolor{background canvas}{bg=violet!20}\fi
    \begin{frame}{TASK for November 26th --}
    
    \mytext{
    \begin{enumerate}
        \item Choose a (small) example project
        \item find a Project Management software
        \item create a free account
        \item create a project
        \item brainstorm the different tasks
        \item enter these tasks
        \item create a Gantt-Diagram
        \item create dependencies between tasks
        \item move the end date of the first task
        \item \textbf{this you demonstrate to the class}
    \end{enumerate}
    
    

    % Considering your PM tool,
    % \begin{itemize}
    %     \item How easy is it to create a WBS?
    %     \item How easy is it to create a list of tasks?
    %     \item What data structure / monitoring does it permit?
    %     \item Is it easy to swap between representations?
    %     \item Is it easy to work with for a group?
    %     \item Does it integrate with other stuff?
    % \end{itemize}
    
    % \framebreak
 
    % Considering your project
    %     \begin{enumerate}
    %         \item What main work packages can you make out?
    %         \item What main milestones can you make out?
    %         \item What responsible persons do you see?
    %     \end{enumerate}

    
    }
\end{frame}
    
\endgroup


% \begin{frame}{Phase Oriented Project Management}
%     \resizebox{!}{0.7\textheight}{
%     \input{tikz/tikz-phase-oriented-project-management}
%     }
    
%     Source: \cite{schiersmannthiel}
% \end{frame}

% \begin{frame}{Typical Main Project Phases (Suchanek)}
%     \resizebox{!}{0.8\textheight}{
%         \input{tikz/tikz-chain-projectphases-details}
%     }

%     \framebreak
    
%     \resizebox{\textwidth}{!}{
%     \input{gantt-mains}
%     }
% \end{frame}



% \begin{frame}{Project Launch per Kerzner}

% \resizebox{!}{0.4\textheight}{
% \parbox{\textwidth}{
%     In the launch phase, customer (= project sponsor), project manager (pm), subject matter experts (sme's) and functional leads (leads) are involved. 
%     \begin{description}
%     \item[Phase 1] - Initialisation
%     \begin{itemize}
%         \item Kickoff-Meeting with customer
%     \end{itemize}
%     \item[Phase 2] - Preliminary planning
%     \begin{itemize}
%         \item PM, SMEs, Leads: Scope Management plan
%         \item Perliminary Scope Statement
%     \end{itemize}
%     \item[Phase 3] - Detailed planning
%     \begin{itemize}
%         \item PM, SMEs, Leads: \textbf{Work Breakdown Structure (WBS)} development
%         \item WBS dictionary
%     \end{itemize}
%     \item[Phase 4] Scope baseline
%     \begin{itemize}
%         \item Finalisation of governing documents\\ ("Lastenheft", product requirement documentation)
%     \end{itemize}
%     \end{description}
    
%     Source: \cite[p. 158]{KerznerHarold2017Pm:a}
% }
% }    
% \end{frame}

% \begin{frame}{The Objectives of a Project}

%     Well defined objectives are prerequisite. BUT
%     \begin{itemize}
%         \item objectives may be subject to change.
%         \item different stake holders may have different views.
%     \end{itemize}
    
    
%     Defined roles and institutions are much less easy to change.
    
% \end{frame}



% \begin{frame}{Project Flow Chart}

    
%     \input{tikz/tikz-projectflowchartone.tex}
    
%     \framebreak


%     \input{tikz/tikz-projectflowcharttwo.tex}

    
%     \framebreak
    
%     \input{tikz/tikz-projectflowchartthree.tex}
    

    
    
% \end{frame}







% \begin{frame}{Dimensions}

%     A project
%     \begin{itemize}
%         \item induces change
%         \item involves people 
%         \item requires resources
%         \item has external effects
%     \end{itemize}

%     A project has
%     \begin{itemize}
%         \item technical dimension (`technisches Angebot')
%         \item business dimension ('kaufmännisches Angebot')
%         \item human dimension (?)
%     \end{itemize}

%     \framebreak

%     Every project has multiple dimensions:
%     \begin{itemize}
%             \item technical
%             \item legal
%             \item economic
%             \item social
%             \item ecological
%             \item IT, data
%             \item ...
%     \end{itemize}
    
%     \framebreak
    
%     The environment could be
%     \begin{itemize}
%         \item hostile or friendly
%         \item stable, or dynamic and reactive
%     \end{itemize}
%     % This can be true for
%     % \begin{itemize}
%     %     \item technology
%     %     \item political / legal situation
%     %     \item economy
%     %     \item market
%     % \end{itemize}
    
%     $\rightarrow$ Different management styles may be required
% \end{frame}





% \input{break}

%\subtitle{Projects, Tasks, Goals and Objectives}


%\frame{\titlepage}

% \begin{frame}{Recap}

%     \begin{itemize}%[<+->]
%         \item A project is unique
%         \item A project delivers change
%         \item A project ventures into the unknown
%         \item A project is supposed to deliver value
%         \item A project consists of tasks
%         \item Complex projects have sub-projects
%     \end{itemize}
    
%     % \framebreak
    
%     % \begin{enumerate}%[<+->]
%     %     \item Who wants the project to happen?
%     %     \item How will success be measured?
%     %     \item Where is the project located?
%     %     \item What is required?
%     %     \item How do we get it done?
%     %     \begin{itemize}
%     %         \item Personel
%     %         \item Material
%     %         \item Other stakeholders
%     %         \item $\dots$
%     %     \end{itemize}
%     % \end{enumerate}
% \end{frame}




% \begin{frame}{WBS Work Breakdown Structure}  % depends on: forest
%     \begin{forest}
%     for tree={
%         font=\footnotesize,
%         grow = east,
%         parent anchor = east,
%         child anchor = west,
%         forked edge,
%         inner sep=2pt
%             }
%     [Travel to Jerusalem
%         [arrive]
%         [start walking [face problems]]
%         [wake up]
%     ]
%     \end{forest}

    
% \end{frame}


% \begin{frame}{WBS Work Breakdown Structure}  % depends on: forest
%     \begin{forest}
%     for tree={
%         font=\footnotesize,
%         grow = east,
%         parent anchor = east,
%         child anchor = west,
%         forked edge,
%         inner sep=2pt
%             }
%     [Travel to Jerusalem
%         [Reviewing [create picture book][review]]
%         [Execution [go to airport][pack stuff]]
%         [Preparation [...][book hotel][book flight][check finances / start saving]]
%     ]
%     \end{forest}
    
        
%     \framebreak
    
%     \smartdiagram[sequence diagram]{preparation, execution, reviewing}
    
%     `We cross the bridge once we come to it' - philosophy
    
%     Foresight?
    
% \end{frame}


\begin{frame}{Project Loft Bed }
    \showmybitmap{../shared/img/hochbett2.jpg}
    \mycomment{
  % depends on: forest
\label{sl:loftbed}

        
    
    }
\end{frame}

\begin{frame}{Project Loft Bed }
        \showmybitmap{../shared/img/hochbett0.jpg}
    \end{frame}

\begin{frame}{Project Loft Bed }    
        \showmybitmap{../shared/img/hochbett1.jpg}
    \end{frame}

\begin{frame}{Project Loft Bed }    
        \showmybitmap{../shared/img/hochbett3.jpg}
    \end{frame}

\begin{frame}{Project Loft Bed }    


    \begin{forest}
    for tree={
        font=\footnotesize,
        grow = east,
        parent anchor = east,
        child anchor = west,
        forked edge,
        inner sep=0.5pt
            }
    [Build a loft bed    ]
    \end{forest}
    \end{frame}

\begin{frame}{Project Loft Bed }
    \begin{forest}
    for tree={
        font=\footnotesize,
        grow = east,
        parent anchor = east,
        child anchor = west,
        forked edge,
        inner sep=0.5pt
            }
    [Build a loft bed
        [Execution]
        [Preparation]
    ]
    \end{forest}
    \end{frame}

\begin{frame}{Project Loft Bed }
    \begin{forest}
    for tree={
        font=\footnotesize,
        grow = east,
        parent anchor = east,
        child anchor = west,
        forked edge,
        inner sep=0.5pt
            }
    [Build a loft bed
        [Execution]
        [Preparation [find helpers][order material][organise tools and workshop][finance][make a plan]]
    ]
    \end{forest}
    \end{frame}

\begin{frame}{Project Loft Bed }
    \begin{forest}
    for tree={
        font=\footnotesize,
        grow = east,
        parent anchor = east,
        child anchor = west,
        forked edge,
        inner sep=0.5pt
            }
    [Build a loft bed
        [Execution [assemble][transport][paint or conserve][shape materials][rent space and tools][order material]]
        [Preparation [find helpers][order material][organise tools and workshop][finance][make a plan]]
    ]
    \end{forest}
    \end{frame}

\begin{frame}{Project Loft Bed }

    \begin{forest}
    for tree={
        font=\footnotesize,
        grow = east,
        parent anchor = east,
        child anchor = west,
        forked edge,
        inner sep=0.5pt
            }
    [Build a loft bed
        [Execution [assemble][transport][paint or conserve][shape materials][rent space and tools][order material]]
        [Preparation [find helpers][order material][organise tools and workshop][finance][make a plan [calculate][draw][measure][consult]]]
    ]
    \end{forest}
    
\end{frame}




\begin{frame}{From WBS to Gantt-Diagram}
    
    \showmytikz{../shared/tikz/tikz-PM-smartd-preparation-execution}
    \end{frame}

\begin{frame}{From WBS to Gantt-Diagram}    
    \showmytikz{../shared/tikz/tikz-PM-wbs-to-gantt-sequence}
    \end{frame}

\begin{frame}{From WBS to Gantt-Diagram}    
    \showmytikz{../shared/tikz/tikz-PM-gantt-loftbed-simple}
\end{frame}

\begin{frame}{From WBS to Gantt-Diagram}
    \showmytikz{../shared/tikz/tikz-PM-gantt-loftbed-nolink}

    \showmytikz{../shared/tikz/tikz-PM-gantt-loftbed-linked}
\end{frame}

\begin{frame}{From WBS to Gantt-Diagram}
    
    \mytext{
    
    \begin{itemize}
        \item Every task should end in a milestone
        \item Define dependencies!
        \item Estimate time (carefully)
        \item Build in buffers
        \item Calculate resources
        \begin{itemize}
            \item no idle time
            \item no overlaps
        \end{itemize}
        \item Use visualisation / software
        % \item[] $\rightarrow$ \textbf{reasearch!}
    \end{itemize}
    
    }
\end{frame}


% \begin{frame}{WBS Hold a Conference}  % depends on: forest
%     \begin{forest}
%     for tree={
%         font=\footnotesize,
%         grow = east,
%         parent anchor = east,
%         child anchor = west,
%         forked edge,
%         inner sep=0.5pt
%             }
%     [Hold a conference
%         [Review]
%         [Execution [venue management][stage service][moderation][welcome participants] ]
%         [Physical planning [contracts][bookings] ]
%         [Content planning [curate: choose presenters][Invite][Book space][Book services]]
%         [Preparation [set up team][finance][plan][advertise][conceptualise]]
%     ]
%     \end{forest}
    
% \end{frame}



% \begin{frame}{WBS Expedition}  % depends on: forest
%     \begin{forest}
%     for tree={
%         font=\footnotesize,
%         grow = east,
%         parent anchor = east,
%         child anchor = west,
%         forked edge,
%         inner sep=0.5pt
%             }
%     [Expedition to Antarctica
%         [Review 
%             [publication]
%             [stakeholder report]
%             [closing]
%             ]
%         [Execution  
%             [scientific research]
%             [watch out weather etc.]
%             [navigate]
%             [set sail / logistics]
%         ]
%         [Preparation 
%             [schedule]
%             [formulate goal]
%             [plan expenses]
%             [find business case] 
%             [identify sponsors]
%             [find vessel]
%             [build up team]
%         ]
%     ]
%     \end{forest}
% \end{frame}




% \begin{frame}{Product Breakdown}
%     \begin{forest}
%     for tree={
%         font=\footnotesize,
%         grow = west,
%         parent anchor = west,
%         child anchor = east,
%         forked edge,
%         inner sep=0.5pt
%             }
%     [New Factory
%         [Managment documents
%             [Business case]
%             [Project governance report]
%             [Project closure report]
%         ]
%         [Legal documents
%             [Approved insurance documents]
%             [Usage certificate]
%         ]
%         [Factory requirements
%             [Use requirements document]
%             [Solution design document]
%         ]
%         [Operational factory
%             [Selected factory]
%             [Fitted internal walls]
%             [Installed furniture and fixtures]
%             [Installed IT and office equipment]
%             [Fitted factory machinery]
%             [Marketed factory]
%         ]
%         [Factory Employees]
%         [Launched Factory]
%         ]
%     \end{forest}
    
%     Source: \cite[p. 168]{Roberts}

% \end{frame}

\begin{frame}{Tasks, Work Packages, Milestones}
    
    \mytext{
    \begin{itemize}
        \item Projects are defined by outcomes
        \item Tasks / Work Packages too
        \item Outcomes are \textbf{milestones}
        \item Milestones are \textbf{interface}
    \end{itemize}
    }
\end{frame}

\begin{frame}{Project Cabin}

    \showmybitmap{../shared/img/fitting-stairs2.JPG}
\end{frame}

\begin{frame}{Project Cabin}
    \showmybitmap{../shared/img/fitting-stairs.JPG}
\end{frame}

\begin{frame}{Project Cabin}
    \showmytikz{../shared/tikz/tikz-PM-gantt-cabin-v1}
    \mycomment{
    
    
    1st plan: only actual work considerd
    
    }
\end{frame}

\begin{frame}{Project Cabin}
    \showmytikz{../shared/tikz/tikz-PM-gantt-cabin-v2}
    \mycomment{
    
    
    Procurement takes longer than expected
    
    }
\end{frame}

\begin{frame}{Project Cabin}
    \showmytikz{../shared/tikz/tikz-PM-gantt-cabin-v3}
    \mycomment{
    
    

    One trade waits for the other
    
    }
\end{frame}

\begin{frame}{Project Cabin}
    \showmytikz{../shared/tikz/tikz-PM-gantt-cabin-v4}
    \mycomment{
    
    
    
    3 trades work in month 5. But when is the planning done?
    
    }
\end{frame}

\begin{frame}{Project Cabin}
    \showmytikz{../shared/tikz/tikz-PM-gantt-cabin-v5}
    \end{frame}

\begin{frame}{Project Cabin}
    
    \mytext{
    
    \begin{itemize}
        \item Planning is super important
        \item Planning itself can and must be planned
        \item Plans may change!
        \begin{itemize}
            \item Make flexible plans!
            \item Plan for alternatives!
        \end{itemize}
        \item Planning looks at \textbf{interfaces}
        \item Planning is communication
    \end{itemize}
    
    }
\end{frame}

 

\begin{frame}{Goal Setting / Defining Milestones}
    
    \mytext{
    \begin{itemize}
        \item How to set clear goals
        \item Useful / measurable mile stones
        \item Measure quality, quantity, \textbf{time}
    \end{itemize}
    }
\end{frame}


% \begin{frame}{Product Flow Chart (part)}
%     \tikzstyle{decision} = [diamond, draw, fill=blue!20, 
%         text width=4.5em, text badly centered, node distance=2cm, inner sep=0pt]
%     \tikzstyle{block} = [rectangle, draw, fill=blue!20, 
%         text width=5em, text centered, rounded corners, minimum height=4em]
%     \tikzstyle{line} = [draw, -latex']
%     \tikzstyle{cloud} = [draw, ellipse,fill=red!20, node distance=3cm,
%         minimum height=2em]
        
%     \begin{tikzpicture}[scale=0.2, node distance = 2.2cm, auto]
%     \footnotesize
%         % Place nodes
%         \node [block] (case) {business case};
%         \node [block, below of=case] (gov) {governance document};
%         \node [block, below left of=gov] (use) {user requirements document};
%         \node [block, below right of=gov] (sol) {solution design document};
%         \node [block, below right of=use] (sel) {selected factory};
%         % Draw edges
%         \path [line] (case) -- (gov);
%         \path [line] (gov) -- (use);
%         \path [line] (gov) -- (sol);
%         \path [line] (use) -- (sel);
%         \path [line] (sol) -- (sel);
%     \end{tikzpicture}
    
%     Source: \cite[p. 169]{Roberts}
    
% \end{frame}

\begin{frame}{The Art of Setting Clear Goals}
    
    \mytext{
    \begin{itemize}
        \item Define precisely what you want
        \item ...and from whom
        \item ...and when
        \item ...and how you will check
        \item ...and what happens if...
    \end{itemize}
    }
\end{frame}

\begin{frame}{The Three Constraints}
    \showmytikz[caption={The Three Constraints (Planning)},label=fig:threeconstraints-planning]{../shared/tikz/tikz-threeconstraints}
    \mycomment{
    
    
    Probably the most important slide.
    }
\end{frame}

\begin{frame}{The Three Constraints}
    
    \mytext{
    You can have it \textbf{cheap,} \textbf{fast, }or \textbf{good.}
    
    Choose two!
    
    \vfill
    \begin{itemize}
        \item Which is more `flexible': quality, time, costs?
        \item Is any one factor critical / not disputable?
        \item Do all stake holder see this in the same way?
        \item How do they interact?
    \end{itemize}
    }
\end{frame}


\begin{frame}{SMART}
    
    \mytext{
    \begin{description}
        \item[S] specific (simple, sensible, significant)
        \item[M] measurable (meaningful, motivating)
        \item[A] achievable (agreed, attainable)
        \item[R] relevant (reasonable, realistic and resourced, results-based)
        \item[T] time bound (time-based, time limited, time/cost limited, timely, time-sensitive)
    \end{description}
    
    \vfill
    
    This applies for tasks and for projects.  
    
    German `operational definiert'
    
    }
\end{frame}







\begin{frame}{Locke and Latham}
    
    \mytext{
    Five goal-setting principles for success:
    
    \begin{enumerate}
        \item Clarity
        \item Challenge
        \item Commitment
        \item Feedback
        \item Task complexity
    \end{enumerate}
    
    \cite{LOCKE1968157}
        
    }
\end{frame}

\begin{frame}{Tasks, Milestones}
    
    \mytext{
    \begin{itemize}
        \item s.m.a.r.t. goals can be milestones
        \item tasks need planning also... or do they?
        \item task/milestone definition: carefully!
        \item is the milestone flexible in time, scope, costs?
        \item is there a way to review goals?
        \item whom is the responsibility attributed to?
    \end{itemize}
    }
\end{frame}




% \begin{frame}{Products vs.~Tasks and Time}
%     \begin{itemize}
%         \item Producs need to be produced
%         \item Production must be planned
%         \item Production takes time: tasks take time!
%         \item Move from a product view to a process view
%         \item Time is the most important \emph{influenceable }cost factor
%         \item How can we estimate time?
%     \end{itemize}
% \end{frame}

% \begin{frame}{Scope Creep}
%     \begin{itemize}
%         \item Sponsor realises he wants more
%         \item Management realises product needs more
%         \item Milking the sponsor or milking the mananger?
%         \item What is necessary?
%         \item How can this all be paid for?
%     \end{itemize}
% \end{frame}

% % \begin{frame}{Product Description}
% %     \begin{itemize}
% %         \item Title
% %         \item Purpose
% %         \item Composition
% %         \item Derivation
% %         \item Format
% %         \item Audience
% %         \item Quality criteria
% %         \item Quality method
% %     \end{itemize}
% % \end{frame}

% % \begin{frame}{The Three Constraints}
% %     \begin{tikzpicture}[scale=0.8]
% %         \coordinate (A) at (0,0);
% %         \coordinate (B) at (-3,-5);
% %         \coordinate (C) at (3,-5);
% %         \coordinate (X) at (0,-3.35);
% %         \draw (A) -- (B) node[midway,above,rotate=60] {Costs};
% %         \draw (B) -- (C) node[midway,below] {Quality/Performance/Technology};
% %         \draw (C) -- (A) node[midway,above,rotate=-60] {Time};
% %         \path[draw=blue, -stealth, rotate=90,
% %       postaction={decorate,decoration={text effects along path,
% %         text={\ Within good customer relations \ }, text align=center, 
% %         text effects/.cd, 
% %           text along path, 
% %           every character/.style={fill=white, yshift=-0.5ex}}}]
% %       (X) circle (4);
% %     \end{tikzpicture}
    
% %     Probably the most important slide.
% % \end{frame}

% \begin{frame}{Scope Freeze Milestones}
%     \begin{itemize}
%         \item A date where the scope cannot be altered
%         \item Typically, specifications are set
%         \item Also design freeze, customer approval freeze...
%         \item De-couple sponsor from product
%         \item De-couple project management from product?
%         \item There is agile!!
%     \end{itemize}
% \end{frame}

% \begin{frame}{MoScCoW - Scheme}
%     \begin{description}
%         \item[M] must have: non-negotiable
%         \item[S] should have: not vital, but adds significant value
%         \item[C] could have: nice to have but small impact when left out
%         \item[W] will not have: stuff without priority in the given time frame.
%     \end{description}    
% \end{frame}

\begin{frame}{Plan planning!}
    
    \mytext{
    Planning 
    \begin{itemize}
        \item creates a model or scenarios
        \item gathers information about possible futures
        \item is a process that takes time in itself
        \item may be self-referential and prone to loops
        \item aims at changing the future, creating value
        \item is necessarily simpler than reality
        \item will get more specific the closer the time
        \begin{itemize}
            \item have a very clear plan for tomorrow
            \item another for the near future
            \item another for the far future
            \item and values for orientation
        \end{itemize}
    \end{itemize}
    }
\end{frame}

% \begin{frame}{Famous Words}
% \begin{quote}
%     Culture eats strategy for breakfast.
    

% \end{quote}
%     \medskip
%     \flushright ---Attr.: Peter \noun{Drucker}
    
% \end{frame}



\begingroup
    \ifdefined\setbeamercolor\setbeamercolor{background canvas}{bg=violet!20}\fi
    \begin{frame}{TASK}
    
    \mytext{
    Create a Gantt-Diagram for the Project Loft Bed (slide~\ref{sl:loftbed}), defining tasks, milestones etc.
        \begin{enumerate}
            \item How long will the project take?
            \item What are the most important tasks, and why?
            \item Which tasks can be done in parallel?
            \item How many work hours will you need?
            \item Which skill sets will you need?
        \end{enumerate}
    
    }
\end{frame}
\endgroup
 
 
%  \begin{frame}{Your example projects!}
%     In groups:
%     \begin{itemize}
%         \item Present each other one of your example projects
%         \item Who wanted what?
%         \item Was it delivered in time, budget, quality?
%     \end{itemize}
    
%     Share one example!
%     \parbox{8cm}{    
%     \centering
%     \clock[0.9]{11}{45}{5}
%     }
% \end{frame}